\section{Literature review}
\label{sec:literature-review}

To put our research regarding the application of agentic workflows to the realm of digital
twins, we review three main areas of research. First, we evaluate the major developments in
agentic systems and tool usage. Next, the geospatial context of GeoAI and digital twins are
examined in closer detail. Finally, we touch briefly on the relevant aspects of computational
linguistics, paraphrasing, and its relation to prompt robustness.

\subsection{Tool Use}
\label{sec:tool-use}

MCP, created by Anthropic, one of the leading AI companies, was widely adopted and has emerged
as the standard for tool integration with LLMs \citep{anthropic2024}. It compares users'
questions against a list of tool descriptions it can call. This also includes all necessary
input parameters as well as output formats, which allows for chaining of tools too. These
textual definitions need to be present to the LLM at all times and together with instructions
how to behave, the so called system prompt, and the actual chat message they occupy the LLM's
context window.

The context window has to be seen as a limiting resource, it is the LLM's ``short term memory''
and everything it knows about the user's conversation. Unfortunately, as the context grows
beyond a few thousand tokens, LLM accuracy begins to degrade when the correct solution can be
found from the center of the context window instead of at the edges \citep{liu2024}. A typical
MCP call to list tools and insert it into the prompt can reach well into this range of tokens.
Beyond this issue, as the context window grows the cost of inference does as well. AI companies
offer caching, e.g.\ Anthropic keeps a users cache alive for 5~minutes after the last
interaction, but this could already be unused if the next user with the same tool definitions in
its context windows sends another request after the short 5~minute period
\citep{anthropicdocs}.

\citet{gan2025} also pointed out that many tool descriptions decrease the chance of LLMs finding
the correct tools. It is currently discussed how granular these tools should be.
\citet{paschalides2026} suggest that tools with medium granularity show best results. Anthropic
has a similar approach stating user should combine different REST-API functionalities to a new
abstracted topic e.g.\ implement a tool called `shedule\_event' that internally calls
`list\_users', `list\_events' and `create\_events' which correspond with the REST-API instead of
making the agent chain them together individually \citep{anthropic2025a}. \citet{mastouri2026}
found that most MCP implementations mirror their respective REST-API.

With MCP the descriptions of all tools are available in the context window initially when the
service starts. Therefore the LLM can incorporate the tools description when defining tasks to
solve the spatial problem beforehand. Another positive aspect is that since these descriptions
don't change over time and also are always in the same place within the context window they can
be cached, which in some scenarios can cut down token usage.

The RAG approach on the other hand does not have knowledge about the available tools beforehand.
In this case the LLM plans the first step to solve the spatial problem and performs a similarity
search in the vector database. It then gets handed 5~possible tools that are the most similar to
the input and picks the one that fits best based on their descriptions. With this approach
caching is not possible, since the content and order is different every time.

\citet{zhang2026} compared `Retrieve-then-Plan' (RtP) with `Plan-then-Retrieve' (PtR) and found
that when using PtR LLMs tend to just pick a tool when the retrieval showed it to them whereas
LLMs that have knowledge of the tools beforehand tend to reject tools that might not fully solve
the current step's problems. Another issue they found was that when using PtR is that LLMs guess
possible solutions without any tie to what the toolset actually can solve and use wrong
assumptions when looking for tools. A more robust approach could be to summarize the general
possibilities of the toolbox at start.

In order to wrap tools into APIs and make the system expandable and the tools exchangeable. We
focus on \href{https://ogcapi.ogc.org/processes/}{OGC API Processes}. This is an open standard
defined by the Open Geospatial Consortium (OGC) that allows for consistent and predictable
execution of processing tasks by modularizing them into a set of standard units. These are
referred to as building blocks and are envisioned as a way to integrate geospatial data with any
other type of data available through similar open standards. Facilitating data integration and
providing consistent service definitions is important for the development of complex geospatial
solutions. For example, the next section provides a background on geospatial artificial
intelligence as well as complicated real-time modeling, which both rely heavily on seamless
integration of data and services.

\subsection{GeoAI \& Digital Twins}
\label{sec:geoai}

The intersection of AI, big data, and high performance computing is referred to as geospatial
artificial intelligence, or GeoAI \citep{li2020}. Artificial intelligence has long been
considered in the context of geospatial technology but is currently being developed at a rapid
pace in comparison to previous decades. While some credit can be given to new computing methods,
an important aspect is also the accessibility and community based sharing of a wider volume of
geospatial data \citep{janowicz2020}. The machine learning methods that are useful for
geospatial problems goes far beyond LLMs and includes other deep learning methods such as
recurrent neural networks and autoencoders, as well as regression and classification methods
like random forest, support vector machines, among others. Likewise, use cases seem to span the
entire range of geospatial applications, from agriculture to urban planning to environmental
monitoring and beyond. The reason for this is that the main methods of network analysis, object
detection, data classification, segmentation, and change detection are useful and relevant for
geospatial data in general \citep{boutayeb2025}. LLMs are not directly related to these methods
of geospatial analysis, but they can be used as agentic workflow generators in a geospatial
context \citep{bhattaram2025}. When combined with a chat-based user interface, this allows for
the democratization of geospatial work including GeoAI, especially in the application of urban
planning and digital twins \citep{ye2026}.

\subsection{Prompt Robustness \& Paraphrasing}
\label{sec:paraphrasing}

Paraphrases are well studied in the field of linguistics, and computational linguistics in
particular is interested in the methods of generating paraphrases. The challenge lies in creating
a sentence that is superficially different but still keeps the same underlying meaning. The
meaning of the sentence can be understood through semantic representation theory
\citep{abend2017}. Exchanging one word for a synonym can introduce unintentional connotations,
adding modal verbs and restructuring clauses can create implications that were not present in
the original sentence. A rigorous definition is necessary for the purpose of examining and
evaluating paraphrases. Paraphrases can be grouped in general by the methods with which they
rearrange and change sentence structure and vocabulary \citep{bhagat2013}. From these methods, a
comprehensive typology of paraphrase types can be established \citep{vila2014}.

Paraphrase generation can be seen as the other side of paraphrase detection, which is of
interest to those who want to detect plagiarism. Paraphrase typologies can be used in detection
and generation to categorize the exact type of paraphrase in either case \citep{wahle2023}. In
the context of LLMs, paraphrasing has been shown to have a limited effect on performance
\citep{wahle2024}. When generating paraphrases, the resulting sentence must have the same meaning
as the original. Paraphrases have also been collected into large datasets. The extended typology
paraphrase corpus is one such dataset that has been catalogued using existing paraphrase
typology \citep{kovatchev2018}. These corpora are valuable for data analysis and training, and
are especially useful for providing examples for an LLM to understand and create a particular
paraphrase type.
